\documentclass[12pt]{article}

\usepackage{amsmath}
\usepackage{graphicx}
\usepackage{booktabs}
\usepackage[a4paper, margin=1in]{geometry}

\title{Labwork 4: Neural Network}
\author{Nguyen Xuan Vinh}
\date{May 09, 2026}

\begin{document}

\maketitle

\section{Introduction}

Artificial Neural Networks are computational models inspired by biological neural systems. 
A neural network consists of interconnected neurons organized into multiple layers.

In this labwork, a simple feedforward neural network was implemented completely from scratch using Python and object-oriented programming concepts.

The implementation supports:
\begin{itemize}
    \item Configurable network architectures
    \item Random weight initialization
    \item Weight initialization from text files
    \item Feedforward computation
\end{itemize}

The network was tested using the XOR logical function.

\section{Network Architecture}

The network architecture is defined using a text configuration file.

The file \texttt{network.txt} contains:

\begin{verbatim}
3
2
2
1
\end{verbatim}

This configuration represents:
\begin{itemize}
    \item 3 layers
    \item Input layer: 2 neurons
    \item Hidden layer: 2 neurons
    \item Output layer: 1 neuron
\end{itemize}

This structure is commonly used to solve the XOR classification problem.

\section{Object Oriented Design}

The neural network was implemented using object-oriented programming.

Three main classes were designed:

\subsection{Neuron Class}

The \texttt{Neuron} class stores:
\begin{itemize}
    \item Weights
    \item Bias
    \item Feedforward computation
\end{itemize}

Each neuron computes a weighted sum of its inputs and applies the sigmoid activation function.

\subsection{Layer Class}

The \texttt{Layer} class contains a list of neurons.

Responsibilities of this class include:
\begin{itemize}
    \item Managing neurons inside the layer
    \item Computing outputs for the entire layer
\end{itemize}

\subsection{NeuralNetwork Class}

The \texttt{NeuralNetwork} class manages:
\begin{itemize}
    \item All layers
    \item Network initialization
    \item Loading weights from files
    \item Feedforward propagation
\end{itemize}

This object-oriented design improves readability, scalability, and maintainability compared to using only lists or dictionaries.

\section{Feedforward Implementation}

Feedforward propagation computes the output of the network layer by layer.

For each neuron, the weighted sum is calculated as:

\[
z = \sum_{i=1}^{n} w_i x_i + b
\]

where:
\begin{itemize}
    \item \(w_i\) are weights
    \item \(x_i\) are inputs
    \item \(b\) is the bias
\end{itemize}

The sigmoid activation function is then applied:

\[
\sigma(z) = \frac{1}{1 + e^{-z}}
\]

The resulting output becomes the input for the next layer.

The process continues until the output layer is reached.

\section{Weight Initialization}

The implementation supports two initialization methods.

\subsection{Random Initialization}

Weights and biases can be initialized randomly in the interval:

\[
[0,1]
\]

using Python's random number generator.

\subsection{Initialization from Text File}

The network also supports loading predefined weights from a text file.

The file \texttt{weights.txt} contains:

\begin{verbatim}
20 20 -10
-20 -20 30
20 20 -30
\end{verbatim}

These values are commonly used for solving the XOR problem using sigmoid activation functions.

\section{XOR Experiment}

The XOR logical function is a classical problem in neural networks because it is not linearly separable.

The network was tested using four input combinations:

\[
(0,0), (0,1), (1,0), (1,1)
\]

\subsection{Experimental Results}

The program produced the following outputs:

\begin{table}[h]
\centering
\begin{tabular}{ccc}
\toprule
Input & Output & Expected XOR Result \\
\midrule
(0,0) & 0.000045 & 0 \\
(0,1) & 0.999955 & 1 \\
(1,0) & 0.999955 & 1 \\
(1,1) & 0.000045 & 0 \\
\bottomrule
\end{tabular}
\caption{XOR neural network outputs}
\end{table}

The outputs are very close to the expected XOR values.

\section{Discussion}

The experiment demonstrates that a multilayer neural network can solve the XOR problem successfully.

Important observations include:
\begin{itemize}
    \item A single-layer perceptron cannot solve XOR
    \item Hidden layers enable nonlinear decision boundaries
    \item Feedforward propagation correctly computes outputs through multiple layers
\end{itemize}

The object-oriented implementation also improves code organization and scalability.

The neural network can easily be extended to:
\begin{itemize}
    \item More hidden layers
    \item More neurons
    \item Different activation functions
    \item Training algorithms such as backpropagation
\end{itemize}

\section{Conclusion}

In this labwork, a feedforward neural network was successfully implemented completely from scratch using Python.

The implementation includes:
\begin{itemize}
    \item Configurable network architecture
    \item Object-oriented design
    \item Random initialization
    \item Weight loading from files
    \item Feedforward computation
\end{itemize}

The XOR experiment demonstrates that the neural network produces outputs very close to the expected logical XOR behavior.

This labwork provides practical understanding of neural network architecture and feedforward computation, which are fundamental concepts in Deep Learning.

\end{document}